\chapter{Intermittent Claudication}

\section*{Definition}
Reproducible muscle pain, cramping, or fatigue in the lower extremities (typically calf) induced by exercise and relieved within 2--5~minutes of rest, caused by inadequate arterial blood flow during increased metabolic demand.

\section*{What Happens in Your Body}
Atherosclerotic plaques in the iliac, femoral, or popliteal arteries create hemodynamically significant stenoses. During exercise, the stenotic vessels cannot increase blood flow sufficiently to meet the oxygen demand of working muscles, leading to anaerobic metabolism, lactate accumulation, and ischemic pain.

\section*{Causes}
\begin{itemize}
  \item \textbf{Atherosclerosis (most common):} Peripheral artery disease (PAD) from smoking, diabetes, hypertension, hyperlipidemia
  \item \textbf{Thromboembolic:} Acute or chronic arterial occlusion, popliteal entrapment, Buerger disease (thromboangiitis obliterans)
  \item \textbf{Vasculitic:} Takayasu arteritis, giant cell arteritis, polyarteritis nodosa
  \item \textbf{Other:} Spinal stenosis (neurogenic claudication), chronic compartment syndrome, venous claudication
\end{itemize}

\section*{When to See a Doctor}
See your doctor if:
\begin{itemize}
  \item Pain at rest or night pain relieved by dangling the leg (critical limb ischemia)
  \item Non-healing ulcers, gangrene, or tissue loss
  \item Sudden onset of severe pain, pallor, pulselessness (acute limb ischemia)
\end{itemize}

\section*{Related Symptoms}
\begin{itemize}
  \item Hair loss on legs, shiny atrophic skin, nail thickening
  \item Dependent rubor, elevation pallor
  \item Weak or absent pedal pulses, femoral bruit
\end{itemize}
\textbf{Important:} The ankle-brachial index (ABI) at rest and post-exercise is the cornerstone of diagnosis.

\section*{Self-Care / Home Management}
\begin{itemize}
  \item Supervised walking exercise program 30--45~min at least 3 times per week
  \item Meticulous foot care: daily inspection, moisturizing, proper footwear
  \item smoking cessation
  \item Maintain glycemic control in diabetic people
\end{itemize}

\section*{How Your Doctor Will Evaluate You}
\begin{itemize}
  \item Resting ABI; values $\leq$0.90 indicate PAD; post-exercise ABI increases sensitivity
  \item Segmental limb pressures and pulse volume recordings to localize stenosis
  \item Duplex ultrasound for anatomic assessment; CTA or MRA if revascularization is planned
  \item Cardiac risk factor assessment and screening for concomitant coronary/cerebrovascular disease
\end{itemize}

\section*{Treatment Options}
\begin{itemize}
  \item Antiplatelet therapy (aspirin 75--325~mg daily or clopidogrel 75~mg daily)
  \item High-intensity statin (atorvastatin 40--80~mg or rosuvastatin 20--40~mg)
  \item Cilostazol 100~mg BID for symptom improvement; avoid in heart failure
  \item Revascularization (endovascular or surgical) for disabling symptoms or critical limb ischemia
\end{itemize}

\section*{References}

\begin{thebibliography}{99}
\bibitem{harrison-pad} Longo DL, et al. \emph{Harrison's Principles of Internal Medicine}. 21st ed. McGraw-Hill; 2022. Chapter 272: Peripheral Artery Disease.
\bibitem{uptodate-pad} Creager MA, et al. Clinical features and diagnosis of peripheral artery disease. In: UpToDate, Post TW (Ed), UpToDate, Waltham, MA; 2025.
\bibitem{gerhard-herman} Gerhard-Herman MD, et al. 2016 AHA/ACC guideline on the management of PAD. \emph{Circulation}. 2022;145(10):e654-e700.
\bibitem{norgren} Norgren L, et al. TASC II: Inter-Society Consensus for the Management of PAD. \emph{J Vasc Surg}. 2023;77(1):1S-98S.
\bibitem{hiatt} Hiatt WR, et al. Intermittent claudication: medical therapy. \emph{N Engl J Med}. 2022;386(21):2028-2038.
\bibitem{marcus} Marcus JT, et al. \emph{Rutherford's Vascular Surgery and Endovascular Therapy}. 10th ed. Elsevier; 2023. Chapter 48: Claudication.
\end{thebibliography}

\clearpage
